\chapter{Introduction}
\label{ch:introduction}

\section{Background}
\label{sec:background}

Machine learning (ML) has emerged as a cornerstone of modern data-driven computational systems, with impactful applications spanning automated healthcare diagnostics, algorithmic financial risk assessment, natural language comprehension, autonomous navigation, and geospatial ecological monitoring. As reliance on machine learning expands across high-stakes domains, the requirements for experimental reliability, reproducibility, and rigorous methodological standards have become paramount.

A fundamental yet frequently overlooked facet of machine learning experimentation is the profound influence of \textit{stochasticity} (experimental randomness). In conventional ML pipelines, pseudo-random number generators (PRNGs) govern numerous critical operational phases:
\begin{enumerate}[label=(\roman*)]
    \item \textbf{Data Splitting:} The random partitioning of observational samples into training, validation, and testing partitions;
    \item \textbf{Parameter Initialization:} The stochastic initialization of feature splits, internal thresholds, or numerical solver weights;
    \item \textbf{Resampling and Subspacing:} The stochastic selection of bootstrap samples and candidate feature subsets during ensemble generation (e.g., in Random Forests);
    \item \textbf{Cross-Validation Fold Assignment:} The random allocation of data records across $k$ cross-validation partitions.
\end{enumerate}

Each of these stochastic factors can induce variations in reported performance metrics---sometimes by several percentage points---without any modification to the underlying model architecture, hyperparameter configuration, or raw dataset. 

The Forest CoverType benchmark dataset, originally collected by the US Geological Survey (USGS) and US Forest Service (USFS) and made available via the UCI Machine Learning Repository \cite{blackard1999comparative}, provides an ideal benchmark for examining these phenomena. Encompassing 581,012 observations characterized by 54 cartographic and geological features across seven distinct forest cover categories, the dataset features substantial class imbalance, continuous spatial gradients, and sparse binary encodings. Investigating how experimental randomness manifests across diverse model families on this dataset provides direct insight into model sensitivity, algorithmic robustness, and scientific reproducibility.

\section{Problem Statement}
\label{sec:problem_statement}

In machine learning research and engineering practice, reported model performance is frequently derived from a single arbitrary train-test partition using an unrecorded or default random seed (such as \texttt{random\_state=42} or \texttt{None}). Consequently:
\begin{itemize}
    \item \textbf{Lack of Seed Reporting:} Benchmark performance comparisons often fail to state whether observed gains stem from algorithmic superiority or fortunate seed configurations;
    \item \textbf{Single-Split Volatility:} A single partition provides a high-variance point estimate that fails to reflect model stability under distribution shifts;
    \item \textbf{Differential Sensitivity:} Distinct machine learning paradigms (e.g., decision trees, ensemble bagging, and maximum-margin hyperplanes) exhibit drastically different vulnerabilities to stochastic perturbations;
    \item \textbf{Scientific Non-Reproducibility:} Inability to replicate reported empirical results due to unseeded data shuffles, platform-dependent floating-point rounding, or undocumented preprocessing order.
\end{itemize}

This project addresses the central problem: \textit{How significantly does experimental randomness (data partitioning seeds, fold allocation, and sample size) influence classification performance, metric stability, and reproducibility across distinct machine learning algorithms?}

\section{Motivation}
\label{sec:motivation}

The motivation for this investigation stems from three core imperatives:
\begin{enumerate}
    \item \textbf{The Machine Learning Reproducibility Crisis:} Prominent machine learning literature (e.g., Bouthillier et al. \cite{bouthillier2019unreproducible}, Dodge et al. \cite{dodge2020fine}, Lucic et al. \cite{lucic2018gans}) has demonstrated that seed variance can obscure genuine algorithmic improvements, leading to false claims of algorithmic superiority.
    \item \textbf{Critical Environmental Applications:} In ecological management, such as forest resource inventory and wildfire risk mitigation, unreliable model classifications can lead to detrimental conservation policies and ineffective resource allocation.
    \item \textbf{Sound Educational \& Industry Standards:} Demonstrating the critical necessity of standardized cross-validation, explicit seed logging, and stratified resampling provides essential practical discipline for machine learning students and practitioners.
\end{enumerate}

\section{Objectives}
\label{sec:objectives}

The primary objectives of this mini project are formulated as follows:
\begin{enumerate}
    \item \textbf{Dataset Curation and Exploratory Data Analysis (EDA):} Ingest the 581,012-sample Forest CoverType dataset via \texttt{scikit-learn}, perform comprehensive univariate, bivariate, and correlation analyses, and construct a stratified subsample of 10,000 samples to enable rigorous multi-seed experimental iterations.
    \item \textbf{Baseline Algorithmic Evaluation:} Implement, train, and evaluate three distinct classification paradigms---Decision Tree (CART), Random Forest, and Linear Support Vector Machine (Linear SVM)---on a standardized 80/20 stratified split, recording Accuracy, Macro-Precision, Macro-Recall, and Macro-F1 scores.
    \item \textbf{Split Seed Sensitivity Quantification:} Train each algorithm across five distinct pseudo-random data split seeds ($S \in \{0, 1, 42, 99, 123\}$) while keeping hyperparameters constant, computing mean performance and sample standard deviation ($\sigma$) to quantify stability.
    \item \textbf{Cross-Validation Variability Assessment:} Execute 5-fold stratified cross-validation across three distinct shuffle seeds ($S_{\text{CV}} \in \{11, 22, 33\}$) to benchmark fold-assignment stability against single-split variance.
    \item \textbf{Deterministic Reproducibility Verification:} Conduct back-to-back replicated training and inference runs with fixed seeds to verify bitwise identical prediction output.
    \item \textbf{Empirical Guidelines Formulation:} Synthesize actionable best practices for seed transparency, stratified resampling, and robust experimental reporting.
\end{enumerate}

\section{Scope of the Project}
\label{sec:scope}

The project scope encompasses:
\begin{itemize}
    \item Three foundational supervised classification architectures: single tree (\texttt{DecisionTreeClassifier}), bagging ensemble (\texttt{RandomForestClassifier}), and linear margin classifier (\texttt{LinearSVC}).
    \item The canonical Forest CoverType multiclass dataset (7 ecological classes, 54 attributes).
    \item A stratified 10,000-sample representative subset ensuring identical class proportions to the full population while enabling fast, multi-iteration experimentation.
    \item Robust preprocessing protocols, specifically standard feature scaling applied strictly to continuous variables without data leakage.
    \item Exhaustive evaluation using overall Accuracy, macro-averaged Precision, Recall, and F1-score, alongside comprehensive per-class classification reports.
\end{itemize}

\textit{Exclusions:} Hyperparameter grid searches, deep neural networks, and non-tabular synthetic data generation lie beyond the designated scope of this mini project.

\section{Contributions}
\label{sec:contributions}

The key contributions of this research study are:
\begin{itemize}
    \item \textbf{Controlled Experimental Framework:} Formulated an isolated, modular experimental pipeline in Python that isolates seed stochasticity from hyperparameter variance.
    \item \textbf{Empirical Stability Hierarchy:} Demonstrated that Random Forest exhibits the lowest variance ($\sigma = 0.58\%$) across split seeds, outperforming Decision Trees ($\sigma = 0.81\%$) and Linear SVM ($\sigma = 0.94\%$).
    \item \textbf{Cross-Validation Variance Proof:} Confirmed that stratified 5-fold cross-validation substantially reduces estimator variance compared to single arbitrary holdout splits.
    \item \textbf{Reproducibility Guarantee:} Proved that deterministic seed management in \texttt{scikit-learn} achieves complete bitwise reproducibility across independent execution instances.
\end{itemize}

\section{Organization of the Report}
\label{sec:organization}

The remainder of this report is structured as follows:
\begin{itemize}
    \item \textbf{Chapter \ref{ch:literature_review} (Literature Review):} Reviews seminal literature on machine learning reproducibility, random seed sensitivity, and algorithmic benchmarks on the Forest CoverType dataset.
    \item \textbf{Chapter \ref{ch:methodology} (Proposed Methodology):} Details the mathematical formulations, dataset properties, stratified subsampling strategy, preprocessing protocols, and evaluation metrics.
    \item \textbf{Chapter \ref{ch:system_design} (System Design):} Outlines the end-to-end system architecture, workflow pipelines, and data flow representations.
    \item \textbf{Chapter \ref{ch:implementation} (Implementation):} Describes the software environment, library dependencies, implementation scripts, and modular code architecture.
    \item \textbf{Chapter \ref{ch:results} (Results and Discussion):} Presents tabular benchmarks, visual plots, seed variance distributions, and comparative analyses.
    \item \textbf{Chapter \ref{ch:conclusion} (Conclusion and Future Scope):} Summarizes the primary findings, outlines methodological limitations, and charts prospective research extensions.
    \item \textbf{Appendix \& Logbook:} Contains United Nations Sustainable Development Goals (SDG) alignment, 12-week project review logbook, and comprehensive graphical visual results.
\end{itemize}
